% Hand-written narrative; numbers are macros.
\emph{Insider-threat detection.} Most work uses the CERT data~\cite{glasser2013bridging,lindauer2020cert}
and reports AUC, accuracy or F1~\cite{yuan2021deep,li2021image,tian2025scenarios}. Le and
Zincir-Heywood~\cite{le2021ensembles} evaluate unsupervised ensembles against an investigation
budget expressed as a share of all user sessions over the whole period; we instead fix the number
of alerts \emph{per day}, which is how an analyst queue operates. Facade~\cite{kantchelian2024facade}
shows at industrial scale that judging an action in the context of its user and resource is what
makes precision workable. Tuor et al.\ use LSTMs that learn each user's sequential behavioural
pattern~\cite{yuan2021deep}, while Yuan et al.\ encode user-day feature vectors as images and apply
CNNs~\cite{li2021image}; both achieve strong AUC values but neither evaluates at a fixed daily
alert capacity. Our evaluation instead selects the alert queue directly: the top \B{} scored
user-days are what an analyst sees each day, and every metric is computed on that queue.

\emph{Cloud audit logs.} Isolation Forest~\cite{liu2008isolation} is applied directly to cloud
control-plane logs for credential-abuse detection and evaluated with F1~\cite{haunschild2026devsecops};
cloud anomaly benchmarks target operational faults rather than insiders~\cite{zou2025cloudanobench}.
The audit log structure we use (who, what resource, when) is the same shape as a cloud API log,
and the features we build are causal in the sense that no feature for day $d$ uses any future
event. This property is essential for cloud deployment, where the model must score a user-day as
soon as the daily log closes without knowing what comes later.

\emph{Calibration and explanation.} Conformal $p$-values give distribution-free outlier
tests~\cite{bates2023conformal}. We use them in a pooled setting: rather than maintaining a
calibration set per user (which would be too small for most employees), we pool the calibrated
family scores of all users into a single reference distribution. This is valid because the
per-user $z$-scores are already normalised, so no user dominates the pool.
Security ML often suffers from temporal snooping, inappropriate
metrics and the base-rate fallacy~\cite{arp2022dos}; only a small share of security XAI studies
evaluate explanations with users~\cite{nadeem2023sok}. SHAP~\cite{lundberg2017shap}-based
insider-threat work re-classifies or labels detections
automatically~\cite{grzeczkowicz2025xai,faqihi2026shap}, and alert-fatigue work acts downstream with
supervision~\cite{ndichu2026pact}. In contrast, our explanations are generated at alert time from
the model's own attribution, without a second post-hoc model, and are evaluated with human
reviewers rather than measured only by agreement with a separate explanation method.

\emph{Evaluation methodology.} The field has documented multiple threats to validity in security
ML: threshold-free metrics that do not reflect operational deployment, seed-based confidence
intervals that confound randomness with genuine uncertainty, and tuning on held-out
windows~\cite{arp2022dos}. Our evaluation addresses all three: we report budgeted metrics in
addition to AUC-ROC, compute intervals by resampling insiders and days rather than seeds, and
tune exclusively inside a development period whose last day precedes the first test day.
